% !TeX root = ../regression_logistique.tex
\chapter{Construction de la fonction logistique}
\label{ann:logit}

\orient{%
construire le logit comme composée de la cote et du logarithme ; inverser
l'hypothèse~\eqref{eq:logit} pour obtenir $\sigma$ ; énoncer ce que cette
construction établit et ce qu'elle suppose}{%
exponentielle, logarithme, cote (\annexeref{ann:prerequis}) ;
chapitre~\ref{ch:nuage}}{%
approfondissement ; le fil principal n'emploie que la formule encadrée}

\section{Première transformation : la cote}

Une probabilité $p$ appartient à $[0,1]$. Le rapport
\[
  \frac{p}{1-p}
\]
s'appelle la \emph{cote}. C'est le langage des paris : une cote de $3$ signifie
trois chances contre une. Ce rapport vaut $0$ quand $p=0$, vaut $1$ quand
$p=0{,}5$, et croît sans borne quand $p$ approche $1$.

\begin{table}[htbp]
\centering
\begin{tabular}{rrrr}
\toprule
$p$ & $1-p$ & cote $p/(1-p)$ & $\log$ de la cote \\
\midrule
$0{,}10$ & $0{,}90$ & $0{,}111$ & $-2{,}197$ \\
$0{,}25$ & $0{,}75$ & $0{,}333$ & $-1{,}099$ \\
$0{,}50$ & $0{,}50$ & $1{,}000$ & $0{,}000$ \\
$0{,}75$ & $0{,}25$ & $3{,}000$ & $+1{,}099$ \\
$0{,}90$ & $0{,}10$ & $9{,}000$ & $+2{,}197$ \\
\bottomrule
\end{tabular}
\caption{La cote transporte $[0,1]$ vers $[0,+\infty[$, son logarithme vers
$\R$ tout entier. La dernière colonne est symétrique autour de $p=0{,}5$ : c'est
la propriété que le score doit posséder.}
\label{tab:cotes}
\end{table}

La cote atteint ainsi toutes les valeurs positives ; restent les négatives.

\begin{figure}[htbp]
\centering
\begin{graphique}{Trois axes superposés reliés par des flèches. Le premier porte l'intervalle de 0 à 1 avec 0,5 au milieu. Le deuxième, le demi-axe des cotes de 0 à l'infini avec 1 au repère. Le troisième, la droite réelle entière avec 0 au repère.}
\begin{tikzpicture}[font=\small,
  seg/.style={line width=1.4pt},
  bnd/.style={font=\scriptsize,text=ink!75},
  op/.style={font=\footnotesize,text=accent,align=center},
  arr/.style={-{Latex[length=2.2mm]},thick,ink!55}]

\draw[seg,gryff!70] (0,0) -- (4.4,0);
\node[bnd,anchor=north] at (0,-0.18) {$0$};
\node[bnd,anchor=north] at (4.4,-0.18) {$1$};
\node[font=\footnotesize,text=ink,anchor=south] at (2.2,0.18) {$p$};

\draw[seg,huffle!70] (0,-2) -- (4.4,-2);
\draw[huffle!70,line width=1.4pt,dashed] (4.4,-2) -- (5.1,-2);
\node[bnd,anchor=north] at (0,-2.18) {$0$};
\node[bnd,anchor=north] at (2.05,-2.18) {$1$};
\node[bnd,anchor=north] at (5.2,-2.18) {$+\infty$};
\fill[huffle!70] (2.05,-2.12) rectangle (2.09,-1.88);
\node[font=\footnotesize,text=ink,anchor=south] at (2.2,-1.82)
  {$\dfrac{p}{1-p}$};

\draw[seg,raven!70] (-1.3,-4) -- (5.1,-4);
\draw[raven!70,line width=1.4pt,dashed] (-2,-4) -- (-1.3,-4);
\node[bnd,anchor=north] at (-2.1,-4.18) {$-\infty$};
\node[bnd,anchor=north] at (2.05,-4.18) {$0$};
\node[bnd,anchor=north] at (5.2,-4.18) {$+\infty$};
\fill[raven!70] (2.05,-4.12) rectangle (2.09,-3.88);
\node[font=\footnotesize,text=ink,anchor=south] at (2.2,-3.82)
  {$\log\dfrac{p}{1-p}$};

\draw[arr] (6.4,0) -- (6.4,-1.6);
\node[op,anchor=west] at (6.6,-0.8) {rapport à\\ son complément};
\draw[arr] (6.4,-2.4) -- (6.4,-3.6);
\node[op,anchor=west] at (6.6,-3.0) {logarithme};
\end{tikzpicture}
\end{graphique}
\caption{Les trois intervalles. La cote étire $]0,1[$ sur $]0,+\infty[$ en envoyant
$0{,}5$ sur $1$ ; le logarithme envoie ce demi-axe sur $\R$. La composée est une
bijection croissante de $]0,1[$ sur $\R$, l'une des infinités qui rendent
l'égalité avec le score possible (section~2).}
\label{fig:intervalles}
\end{figure}
\FloatBarrier

\section{Seconde transformation : le logarithme}

Le logarithme envoie $]0,+\infty[$ sur $\R$ : il vaut $-\infty$ près de zéro,
$0$ en $1$, et croît sans borne. En le composant avec la cote, on obtient la
\emph{log-cote}
\[
  \log\frac{p}{1-p},
\]
qui prend toutes les valeurs réelles, comme le score $z$. Les deux quantités
parcourant le même ensemble, poser leur égalité est \emph{possible} :
% \tag ne fait pas avancer le compteur : sans ce pas, l'ancre PDF de cette
% équation entre en collision avec celle de la première équation numérotée.
\begin{equation}
  \boxed{\;\log\frac{p}{1-p}=z=\sum_{j=0}^{n}w_jx_j\;}
  \tag{$\star$}\label{eq:logit}
\end{equation}
\stepcounter{equation}

L'égalité~\eqref{eq:logit} constitue un choix parmi les bijections croissantes
de $]0,1[$ sur $\R$. La réciproque de la fonction de
répartition de la loi normale donne le modèle \emph{probit}, celle de la loi de
Gumbel le modèle \emph{complémentaire log-log} (ch.~\ref{ch:nuage}). Le logit est
donc une hypothèse de modélisation, retenue pour trois raisons : il donne aux
coefficients une lecture directe en log-cotes, il est le lien canonique du modèle
linéaire généralisé binomial, ce qui produit les formules les plus
simples~\cite{mccullagh1989}, et il conduit à une perte convexe
(\annexeref{ann:convexite}). Le modèle en compte deux autres, posées ailleurs :
l'indépendance conditionnelle des observations (\annexeref{ann:cout}) et le choix
des variables retenues.

Les annexes suivantes tirent les conséquences de~\eqref{eq:logit} : l'expression
de $p$ en fonction de $z$ ci-dessous, les propriétés de $\sigma$
(\annexeref{ann:sigmoide}), la perte (\annexeref{ann:cout}) et le gradient
(\annexeref{ann:gradient}).

\section{Expression de \texorpdfstring{$p$}{p} en fonction de \texorpdfstring{$z$}{z}}

L'équation~\eqref{eq:logit} donne $z$ à partir de $p$, alors que le programme
part de $z$ pour obtenir la probabilité. L'inversion consiste à isoler $p$ ;
chaque ligne porte à droite l'opération qui la produit.

\begin{align}
  \log\frac{p}{1-p} &= z \notag\\[2pt]
  e^{\log\frac{p}{1-p}} &= e^{z}
    &&\text{on applique } e^{\;\cdot\;} \text{ aux deux membres} \notag\\[2pt]
  \frac{p}{1-p} &= e^{z}
    &&\text{car } e^{\log a}=a \notag\\[2pt]
\end{align}

\slidebreak

\begin{align}
  p &= e^{z}(1-p)
    &&\text{on multiplie les deux membres par } 1-p \notag\\[2pt]
  p &= e^{z}-e^{z}p
    &&\text{on développe} \notag\\[2pt]
  p+e^{z}p &= e^{z}
    &&\text{on rassemble les } p \notag\\[2pt]
  p\,(1+e^{z}) &= e^{z}
    &&\text{on met } p \text{ en facteur} \notag\\[2pt]
  p &= \frac{e^{z}}{1+e^{z}}
    &&\text{on divise par } 1+e^{z},\ \text{qui n'est jamais nul} \label{eq:sigma-brute}
\end{align}

Cette écriture se met sous une forme plus courante en divisant numérateur et
dénominateur par $e^{z}$ :
\begin{align}
  p &= \frac{e^{z}}{1+e^{z}}
     = \frac{e^{z}/e^{z}}{(1+e^{z})/e^{z}}
    &&\text{on divise en haut et en bas par } e^{z}\neq0 \notag\\[2pt]
    &= \frac{1}{1/e^{z}+1}
    &&\text{car } e^{z}/e^{z}=1 \notag\\[2pt]
    &= \frac{1}{1+e^{-z}}
    &&\text{car } 1/e^{z}=e^{-z} \label{eq:sigma}
\end{align}

\begin{resultat}[Fonction logistique]
\[
  \sigma(z)=\frac{1}{1+e^{-z}}=\frac{e^{z}}{1+e^{z}} .
\]
Les deux écritures désignent la même fonction, obtenue en inversant
l'hypothèse~\eqref{eq:logit}.
\end{resultat}

Chacune des deux formes échoue là où l'autre tient. Pour $z$ très négatif,
$e^{-z}$ devient
gigantesque et $1/(1+e^{-z})$ déborde la capacité de la machine, tandis que
$e^{z}/(1+e^{z})$ rend $0$. Pour $z$ très positif, c'est
l'inverse. Le programme utilise donc l'une ou l'autre selon le signe de $z$
(\annexeref{ann:stabilite}).

\section{Vérification numérique}

La table~\ref{tab:cotes} se relit dans l'autre sens. Pour $z=1{,}099$ :
\[
  e^{-1{,}099}=0{,}3333,\qquad
  \sigma(1{,}099)=\frac{1}{1+0{,}3333}=\frac{1}{1{,}3333}=0{,}75 .
\]
On retrouve la ligne $p=0{,}75$, dont la log-cote valait $+1{,}099$. Pour
$z=0$ :
\[
  \sigma(0)=\frac{1}{1+e^{0}}=\frac{1}{1+1}=0{,}5 .
\]
Un élève situé exactement sur la droite reçoit une probabilité de $0{,}5$ : la
frontière (ch.~\ref{ch:nuage}) est la ligne de niveau $0{,}5$ du modèle.

\begin{figure}[htbp]
\centering
\begin{graphique}{Courbe de la fonction logistique, en S, croissante de 0 vers 1, valant 0,5 à l'origine, avec ses deux asymptotes horizontales en pointillé.}
\begin{tikzpicture}
\begin{axis}[courseplot,height=6.4cm,axis lines=middle,
  xlabel={score $z$},ylabel={$\sigma(z)$},
  xmin=-6.5,xmax=6.5,ymin=-.12,ymax=1.15,ytick={0,.25,.5,.75,1},
  samples=200,domain=-6.5:6.5]
  \addplot[gridgray,thick,dashed,domain=-6.5:6.5]{1};
  \addplot[accent,very thick]{1/(1+exp(-x))};
  \addplot[only marks,mark=*,warning,mark size=2.6pt]
    coordinates {(0,0.5) (1.099,0.75) (-1.099,0.25)};
  \node[warning,anchor=west,font=\small] at (axis cs:1.3,.78){$(1{,}099\,;\,0{,}75)$};
  \node[warning,anchor=east,font=\small] at (axis cs:-1.3,.22){$(-1{,}099\,;\,0{,}25)$};
  \node[ink,anchor=south west,font=\small] at (axis cs:.15,.5){$\sigma(0)=0{,}5$};
\end{axis}
\end{tikzpicture}
\end{graphique}
\caption{La fonction logistique envoie $\R$ dans l'intervalle $]0,1[$, sans
jamais atteindre les bornes, et passe par $0{,}5$ à l'origine.}
\label{fig:sigmoide}
\end{figure}
\FloatBarrier

\begin{vigilance}[Valeurs prises par $\sigma$]
Pour $z$ fini, $1+e^{-z}$ reste strictement plus grand que $1$, donc $\sigma(z)$
reste strictement compris entre $0$ et $1$. Toute probabilité produite par le
modèle appartient donc à l'intervalle ouvert $]0,1[$. En machine, $\sigma(40)$
vaut pourtant exactement $1$ après arrondi : l'écart disparaît à partir d'un
score d'environ $37$, avec les conséquences que cela entraîne dans un logarithme
(\annexeref{ann:stabilite}).
\end{vigilance}
